%==============================================================================
%  TEMPLATE LAPORAN SYSTEMATIC LITERATURE REVIEW (SLR)
%  Program Diploma IV Teknik Informatika
%  Universitas Logistik dan Bisnis Internasional (ULBI)
%
%  Hasil konversi LaTeX dari "templateslr.docx".
%
%  Kompilasi:
%     pdflatex main
%     bibtex   main
%     pdflatex main
%     pdflatex main
%  atau cukup:  latexmk -pdf main.tex
%==============================================================================
\documentclass[12pt,a4paper,oneside]{report}

%------------------------------------------------------------------ encoding & bahasa
\usepackage[T1]{fontenc}
\usepackage[utf8]{inputenc}
\usepackage[english,bahasa]{babel}      % bahasa utama: Indonesia

%------------------------------------------------------------------ font Times New Roman
\usepackage{newtxtext,newtxmath}        % pengganti Times New Roman

%------------------------------------------------------------------ tata letak halaman
% Margin mengikuti templateslr.docx (kiri/kanan 3 cm). Untuk format jilid
% Indonesia (kiri 4 cm) ubah "left=3cm" menjadi "left=4cm".
\usepackage[a4paper,top=3cm,bottom=3cm,left=3cm,right=3cm]{geometry}
\usepackage{setspace}
\onehalfspacing                         % spasi 1,5

%------------------------------------------------------------------ paket pendukung
\usepackage{graphicx}
\usepackage{enumitem}
\usepackage{array}
\usepackage{longtable}
\usepackage{booktabs}
\usepackage{ragged2e}
\usepackage{titlesec}
\usepackage{fancyhdr}
\usepackage{caption}
% hypertexnames=false -> anchor halaman dibuat berurutan secara internal,
% sehingga reset nomor halaman tiap BAB tidak menimbulkan bentrokan anchor.
\usepackage[hidelinks,hypertexnames=false]{hyperref}

%------------------------------------------------------------------ nama bagian (Bahasa Indonesia)
\addto\captionsbahasa{%
  \renewcommand{\contentsname}{DAFTAR ISI}%
  \renewcommand{\listfigurename}{DAFTAR GAMBAR}%
  \renewcommand{\listtablename}{DAFTAR TABEL}%
  \renewcommand{\bibname}{DAFTAR PUSTAKA}%
  \renewcommand{\chaptername}{BAB}%
}

%------------------------------------------------------------------ penomoran sub-bab
% Bab heading memakai angka Romawi (BAB I, BAB II), namun penomoran sub-bab,
% gambar, dan tabel memakai angka Arab seperti pada dokumen sumber.
\setcounter{secnumdepth}{3}
\renewcommand{\thesection}{\arabic{section}}
\renewcommand{\thesubsection}{\thesection.\arabic{subsection}}
\renewcommand{\thesubsubsection}{\thesubsection.\arabic{subsubsection}}

%------------------------------------------------------------------ format judul BAB
\titleformat{\chapter}[display]
  {\centering\normalfont\bfseries\large}
  {BAB \Roman{chapter}}{6pt}{\MakeUppercase}
\titlespacing*{\chapter}{0pt}{0pt}{20pt}

\titleformat{\section}{\normalfont\bfseries}{\thesection.}{0.6em}{}
\titleformat{\subsection}{\normalfont\bfseries}{\thesubsection}{0.6em}{}
\titleformat{\subsubsection}{\normalfont\bfseries\itshape}{\thesubsubsection}{0.6em}{}

%------------------------------------------------------------------ caption gambar & tabel
\renewcommand{\thefigure}{\arabic{chapter}.\arabic{figure}}
\renewcommand{\thetable}{\arabic{chapter}.\arabic{table}}
\captionsetup[figure]{name=Gambar,labelsep=space,justification=centering,font=bf}
\captionsetup[table]{name=Tabel,labelsep=space,justification=centering,font={bf,it},position=above}

%------------------------------------------------------------------ reset nomor halaman tiap BAB
% Halaman ditulis dengan format <Romawi-Bab>-<halaman>, mis. I-1, II-6.
\makeatletter
\let\slr@old@chapter\@chapter
\def\@chapter[#1]#2{\setcounter{page}{1}\slr@old@chapter[{#1}]{#2}}
\makeatother

\newcommand{\mulaiisi}{%
  \clearpage
  \setcounter{chapter}{0}%
  \renewcommand{\thepage}{\Roman{chapter}-\arabic{page}}%
}

%------------------------------------------------------------------ footer: nomor halaman tengah bawah
\pagestyle{fancy}
\fancyhf{}
\fancyfoot[C]{\thepage}
\renewcommand{\headrulewidth}{0pt}
\fancypagestyle{plain}{%
  \fancyhf{}\fancyfoot[C]{\thepage}\renewcommand{\headrulewidth}{0pt}}

%==============================================================================
%  METADATA  ---  ubah bagian ini sesuai laporan Anda
%==============================================================================
\newcommand{\judulID}{TINJAUAN STUDI LITERATUR PERAN \textit{MACHINE LEARNING} DALAM MEMPREDIKSI \textit{HAPPINESS} DAN KESEJAHTERAAN SOSIAL-EKONOMI}
\newcommand{\judulEN}{A SYSTEMATIC LITERATURE REVIEW ON THE ROLE OF MACHINE LEARNING IN PREDICTING HAPPINESS AND SOCIOECONOMIC WELL-BEING}
\newcommand{\jenisID}{LAPORAN INTERNSHIP 2}
\newcommand{\jenisEN}{INTERNSHIP 2 REPORT}
\newcommand{\subjudulID}{Diajukan Sebagai Salah Satu Syarat Kelulusan Mata Kuliah Internship~2 Pada Program Pendidikan Diploma IV Teknik Informatika.}
\newcommand{\subjudulEN}{\textit{Submitted as one of the graduation requirements for the Internship~2 course in the Informatics Engineering Diploma IV Education Program}}
\newcommand{\penulis}{ROFI NAFIIS ZAIN}
\newcommand{\nim}{1.21.4.017}
\newcommand{\prodi}{PROGRAM DIPLOMA IV TEKNIK INFORMATIKA}
\newcommand{\fakultas}{SCHOOL OF TECHNOLOGY INFORMATION}
\newcommand{\universitas}{UNIVERSITAS LOGISTIK DAN BISNIS INTERNASIONAL}
\newcommand{\kota}{Bandung}
\newcommand{\bulantahun}{Juni 2026}
\newcommand{\logo}{images/logo-ulbi.png}

%==============================================================================
\begin{document}
\pagestyle{empty}

%------------------------------------------------------------------ HALAMAN SAMPUL (Indonesia)
\begin{titlepage}\centering
  \begin{spacing}{1.2}
    {\bfseries\large \judulID\par}
    \vspace{1.5em}
    {\bfseries\large \jenisID\par}
    \vspace{1em}
    {\itshape \subjudulID\par}
  \end{spacing}
  \vspace{1.5em}
  \includegraphics[width=0.45\textwidth]{\logo}\par
  \vspace{1.5em}
  {\bfseries Dibuat oleh,\par}
  \vspace{0.5em}
  {\bfseries \nim\quad \penulis\par}
  \vfill
  {\bfseries \prodi\par}
  {\bfseries \fakultas\par}
  {\bfseries \universitas\par}
\end{titlepage}

%------------------------------------------------------------------ HALAMAN SAMPUL (Inggris)
\begin{titlepage}\centering
  \begin{spacing}{1.2}
    {\bfseries\large \judulEN\par}
    \vspace{1.5em}
    {\bfseries\large \jenisEN\par}
    \vspace{1em}
    {\itshape \subjudulEN\par}
  \end{spacing}
  \vspace{1.5em}
  \includegraphics[width=0.45\textwidth]{\logo}\par
  \vspace{1.5em}
  {\bfseries Created by,\par}
  \vspace{0.5em}
  {\bfseries \nim\quad \penulis\par}
  \vfill
  {\bfseries APPLIED BACHELOR PROGRAM INFORMATICS ENGINEERING\par}
  {\bfseries \fakultas\par}
  {\bfseries \universitas\par}
\end{titlepage}

%------------------------------------------------------------------ LEMBAR PENGESAHAN (hasil scan)
% Sisipkan halaman pengesahan/persetujuan yang sudah ditandatangani di sini,
% misalnya hasil scan dalam bentuk gambar atau PDF:
%   \includegraphics[width=\textwidth]{images/lembar-pengesahan.jpg}
%   \includegraphics[width=\textwidth]{images/lembar-persetujuan.jpg}
%   \includegraphics[width=\textwidth]{images/lembar-orisinalitas.jpg}

%==============================================================================
%  BAGIAN AWAL  (penomoran i, ii, iii, ...)
%==============================================================================
\pagenumbering{roman}
\pagestyle{fancy}

%------------------------------------------------------------------ ABSTRAK
\chapter*{ABSTRAK}
\addcontentsline{toc}{chapter}{ABSTRAK}
\justifying
Tuliskan abstrak Bahasa Indonesia di sini: latar belakang singkat, tujuan
\textit{Systematic Literature Review} (SLR), metode (mis. PRISMA), ringkasan
temuan, dan kontribusi penelitian. Satu paragraf, umumnya 150--250 kata.

\par\vspace{1em}
\noindent\textbf{Kata Kunci:} \textit{Machine Learning, Well-Being Prediction,
Life Satisfaction, Happiness Index, Socio-economic Factors}

%------------------------------------------------------------------ ABSTRACT
\chapter*{ABSTRACT}
\addcontentsline{toc}{chapter}{\textit{ABSTRACT}}
\justifying
\textit{Write the English abstract here as a faithful translation of the
Indonesian abstract above.}

\par\vspace{1em}
\noindent\textbf{Keyword:} \textit{Machine Learning, Neural Decoding,
Alphanumeric, Brain-Computer Interface, Systematic Literature Review (SLR)}

%------------------------------------------------------------------ KATA PENGANTAR
\chapter*{KATA PENGANTAR}
\addcontentsline{toc}{chapter}{KATA PENGANTAR}
\justifying
Puji syukur penulis panjatkan kepada Tuhan Yang Maha Esa, karena atas berkat
Rahmat-Nya, laporan dengan judul \textit{``\judulEN''} dapat diselesaikan tepat
pada waktunya. Pada kesempatan ini penulis menyampaikan terima kasih kepada:

\begin{enumerate}[label=\arabic*., leftmargin=*, itemsep=2pt]
  \item Kedua orang tua penulis atas dukungan, semangat, dan doanya.
  \item Rektor \universitas.
  \item Ketua Program Studi Sarjana Terapan Teknik Informatika.
  \item Dosen pembimbing utama atas bimbingan dan masukannya.
  \item Dosen pembimbing kedua atas bimbingan dan masukannya.
  \item Dosen penguji atas masukan dan dukungannya.
  \item Teman-teman yang telah membantu dan memberi dukungan.
\end{enumerate}

Penulis menyadari laporan ini masih memiliki keterbatasan, sehingga kritik dan
saran dari pembaca sangat diharapkan.

\par\vspace{2em}
\hfill\begin{tabular}{c}\kota, \bulantahun\\[3em]\penulis\end{tabular}

%------------------------------------------------------------------ DAFTAR ISI / GAMBAR / TABEL / LAMPIRAN
\tableofcontents
\listoffigures
\listoftables

\chapter*{DAFTAR LAMPIRAN}
\addcontentsline{toc}{chapter}{DAFTAR LAMPIRAN}
\begin{description}[leftmargin=0pt,labelindent=0pt]
  \item[] Lampiran 1\quad Formulir Kemajuan Bimbingan Internship 2
\end{description}

%==============================================================================
%  BAGIAN ISI  (penomoran I-1, II-1, ...)
%==============================================================================
\mulaiisi

%------------------------------------------------------------------ BAB I
\chapter{PENDAHULUAN}\justifying

\section{Latar Belakang}
Uraikan latar belakang penelitian: fenomena/masalah, urgensi, ringkasan
penelitian terdahulu, dan posisi penelitian ini. Sitasi memakai gaya IEEE,
contoh \cite{kaufman2022,ruggeri2020}. Beberapa sumber sekaligus
\cite{casaccia2020,izumi2021}.

\section{Identifikasi Masalah}
Adapun identifikasi masalah dari penelitian ini sebagai berikut.
\begin{enumerate}[label=\arabic*., leftmargin=*, itemsep=2pt]
  \item Bagaimana integrasi teknik \textit{machine learning} dapat meningkatkan
        prediksi kesejahteraan dan kepuasan hidup di berbagai konteks sosioekonomi?
  \item Apa peran PDB dan ketimpangan pendapatan dalam memengaruhi kebahagiaan
        dan kualitas hidup, dan bagaimana \textit{machine learning} dapat
        memberikan wawasan yang lebih mendalam tentang hubungan ini?
  \item Apa saja aplikasi baru dari indeks kebahagiaan dalam mengukur kemajuan
        masyarakat di luar PDB, dan bagaimana indeks-indeks ini dapat
        ditingkatkan menggunakan model prediktif?
\end{enumerate}

\section{Tujuan dan Manfaat}
\subsection{Tujuan Penelitian}
Adapun tujuan dari penelitian ini sebagai berikut.
\begin{enumerate}[label=\arabic*., leftmargin=*, itemsep=2pt]
  \item Menganalisis kontribusi \textit{machine learning} terhadap akurasi
        prediksi kesejahteraan.
  \item Mengungkap hubungan antara PDB, ketimpangan pendapatan, dan kebahagiaan.
  \item Mengevaluasi efektivitas indeks kebahagiaan yang ada.
\end{enumerate}

\subsection{Manfaat Penelitian}
Adapun manfaat penelitian ini sebagai berikut.
\begin{enumerate}[label=\arabic*., leftmargin=*, itemsep=2pt]
  \item Menyediakan metode prediksi kesejahteraan yang lebih akurat.
  \item Memberikan pemahaman hubungan PDB, ketimpangan pendapatan, dan kebahagiaan.
  \item Menghasilkan pengembangan indeks kebahagiaan berbasis model prediktif.
\end{enumerate}

\section{Ruang Lingkup}
Adapun ruang lingkup dari penelitian ini adalah.
\begin{enumerate}[label=\arabic*., leftmargin=*, itemsep=2pt]
  \item Menggunakan \textit{literature} dari \textit{database} Scopus (Q1--Q4).
  \item Menerapkan \textit{filter} tahun publikasi 2019--2024.
  \item Menggunakan \textit{keyword} spesifik untuk pencarian literatur.
  \item \textit{Tools} yang digunakan adalah WATASE UAKE dengan metode utama PRISMA.
\end{enumerate}

%------------------------------------------------------------------ BAB II
\chapter{LANDASAN TEORI}\justifying

\section{State of the Art}
Uraikan posisi penelitian terhadap karya terdahulu.

\section{Hasil Survey Studi Literatur}
Bagian ini menyajikan hasil survei literatur berdasarkan pendekatan PRISMA.

% Contoh tabel panjang (dapat melintasi beberapa halaman) -- Tinjauan Pustaka.
\begin{longtable}{>{\raggedright}p{1.6cm} p{2.2cm} p{2.2cm} p{4cm} p{3cm}}
  \caption{Tinjauan Pustaka}\label{tab:tinjauan}\\
  \toprule
  \textbf{Penulis} & \textbf{Metode} & \textbf{Model} & \textbf{Tujuan} & \textbf{Evaluasi Metrik}\\
  \midrule
  \endfirsthead
  \multicolumn{5}{l}{\small\itshape Lanjutan Tabel \thetable}\\
  \toprule
  \textbf{Penulis} & \textbf{Metode} & \textbf{Model} & \textbf{Tujuan} & \textbf{Evaluasi Metrik}\\
  \midrule
  \endhead
  \midrule \multicolumn{5}{r}{\small\itshape bersambung ke halaman berikutnya}\\
  \endfoot
  \bottomrule
  \endlastfoot
  \cite{kaufman2022} & Regresi & Lasso Regression & Memprediksi kebahagiaan lansia & RMSE = 1.195; R\textsuperscript{2} = 0.319\\
  \cite{ruggeri2020} & Regresi & Ridge Regression & Estimasi GDP subnasional & R\textsuperscript{2} = 0.95; MSE = 0.01\\
  \cite{casaccia2020} & Klasifikasi & Random Forest & Prediksi kesejahteraan mental & Akurasi = 0.921; AUC = 0.966\\
\end{longtable}

\section{Tinjauan Waktu dan Wilayah Studi}
% Contoh penyisipan gambar.
\begin{figure}[ht]\centering
  \includegraphics[width=0.5\textwidth]{\logo}
  \caption{Tahun publikasi studi}\label{fig:tahun}
\end{figure}
Sebagaimana ditunjukkan pada Gambar~\ref{fig:tahun}, \dots

\section{Area Studi dan Tren Metodologi Terdahulu}
\section{Analisis Word Cloud}
\section{Distribusi Hipotesis dan Variabel}

\section{\textit{Watase Network Analysis} (WNA)}
\subsection{Analisis Teori}
\subsection{Analisis Konteks}

\section{Gap Penelitian}
\subsection{Evidence Gap}
\subsection{Empirical Gap}

\section{Analisis Pendekatan}

\section{Analisis Konten}
\subsection{Integrasi \textit{machine learning} untuk prediksi kesejahteraan (RQ1)}
\subsection{Peran PDB dan ketimpangan pendapatan (RQ2)}
\subsection{Aplikasi baru indeks kebahagiaan (RQ3)}

%------------------------------------------------------------------ BAB III
\chapter{METODOLOGI PENELITIAN}\justifying

\section{Diagram Alur Penelitian}
\section{Penjelasan Diagram Alur Metodologi Penelitian}
\subsection{\textit{Identification}}
\subsection{\textit{Screening}}
\subsection{\textit{Included}}

%------------------------------------------------------------------ BAB IV
\chapter{PENUTUP}\justifying

\section{Kesimpulan dan Saran}
\subsection{Kesimpulan}
\subsection{Saran}

\section{Lampiran--Lampiran}
% \textit{Lampiran 1 Formulir Kemajuan Bimbingan Internship 2}
% \includegraphics[width=0.8\textwidth]{images/lampiran1.jpg}

%==============================================================================
%  DAFTAR PUSTAKA  (gaya IEEE)
%==============================================================================
\renewcommand{\thepage}{\arabic{page}}   % daftar pustaka memakai nomor biasa
\addcontentsline{toc}{chapter}{DAFTAR PUSTAKA}
\bibliographystyle{IEEEtran}
\bibliography{references}

\end{document}
